% Nouveau EGA V (cote 25): archive cote 25, site page 2 (PDF page 3).
% typescript. Draft compared: none.
% AI-transcribed by Claude Opus 5.5 (claude-opus-5-5, adaptive thinking, effort high), Message Batches API, batch pilot (msgbatch_0189Nhf1xB615YWnztubtLfc), 2026-11-12; not yet checked by eye.
\scanpage{25/2}
\margin{\unsure{Nouveau EGA V}} \mspage{0}

EGA IV 20 \quad \uline{Sections hyperplanes et projections coniques.} (Etat \unsure{000|0})

1. Notations.

2. Section hyperplane générique: propriétés locales.

3. Section hyperplane générique: irréductibilité et connexité géo\unsure{métriques.}

4. Section hypérplane variable: sections ``assez g\unsure{énérales|éńérales}''.

5. Théorèmes du\unsure{type|ttype} Seidenberg.

6. Connexion d'une section hyperplane quelconque.

7. Application a des constructions de sections hyperplanes et multisec-%
tions spéciales.

8. Dimension de l'ensem\xout{m}ble  des hyperplans exceptionnels.

9. Changement d'immersion projective.

10. Pinceaux de sections hyperplanes, et fibrations de variétés éclatées

11. Grassmaniennes.

12. Généralisation de résultats antérieurs à des sections linéaires.

13. Morphismes élémentaires et théorème de M. Artin.

14. Projections coniques.

15. Axiomatisation de certains des résultats précédents.

16. Traductions en langage de systèmes linéaires.

\uline{Commentaires}. Cette rédaction donne pêle mêle  une esquisse détaillée
de l'ensemble des résultats qui devraient figurer dans une rédaction
finale. Pour parvenir à cette dernière, il faudra une réorganisation de
fond en comble  du présent Etat 0, la première tâche étant probablement
l'établissement d'un nouveau plan (dans lequel sans doute  les actuels
nºs 11,12,14,15  viendront beaucoup plus tôt). Je n'ai pas encore rédi\unsure{-}
le nº 16 , qui en principe ne doit offrir aucune difficulté ni
influer en aucune manière sur les nºs précédents, puisqu'il s'agit d'u\unsure{n}

\margin{2}
